\chapter{Anexos}

\section{Funciones para Datos Tokenizados}

A continuación, se muestran funciones adicionales del proyecto para con los datos tokenizados:

\subsection{Tokenización de los datos}

\begin{center}
\textbf{Listing A.1.} Función para Tokenizar los Datos
\end{center}

\begin{lstlisting}[style=pythoncode]
def save_tokenized_data(df_dict, max_length=None):
    """
    Convierte los datos tokenizados (diccionario de arrays) a DatasetDict de
    Hugging Face
    y los guarda en disco.
    """

    tokenizer=load_tokenizer()

    hf_datasets = DatasetDict()

    for subset, df in df_dict.items():
        # Tokenizar los comentarios
        tokenized = tokenizer(
            df["comment"].tolist(),
            padding=True,
            truncation=max_length is not None,
            max_length=max_length,
            return_tensors="np"
        )

        # Crear Dataset
        hf_dataset = Dataset.from_dict({
            "input_ids": tokenized["input_ids"],
            "attention_mask": tokenized["attention_mask"],
            "labels": np.array(df["label"])
        })

        hf_datasets[subset] = hf_dataset
\end{lstlisting}

% =========================
% Página 70
% =========================

\vspace{0.5cm}

\begin{lstlisting}[style=pythoncode]
    tokenize_path = Path(os.environ.get("TOKENIZE_PATH"))
    tokenize_path.mkdir(parents=True, exist_ok=True)
    hf_datasets.save_to_disk(tokenize_path)

    return {
        "message": "Tokenización de la data exitosa",
        "ruta_de_guardado": str(tokenize_path)
    }
\end{lstlisting}

\section{Muestra de un Dato Tokenizado}

A continuación, se muestra la función que mostró un ejemplo de dato tokenizado en el capítulo 9, además de un segundo ejemplo de salida.

\begin{center}
\textbf{Listing A.2.} Función para mostrar un Dato Tokenizado
\end{center}

\begin{lstlisting}[style=pythoncode]
def show_tokenized_example(subset="train", idx=None):
    """
    Muestra los datos tokenizados.
    """

    # Cargar el DatasetDict desde disco
    tokenize_path = Path(os.environ.get("TOKENIZE_PATH"))
    hf_datasets = DatasetDict.load_from_disk(tokenize_path)

    dataset = hf_datasets[subset]
    if idx is None:
        idx = random.randint(0, len(dataset) - 1)

    example = dataset[idx]
    tokenizer = load_tokenizer()
    text = tokenizer.decode(example['input_ids'], skip_special_tokens=True)
    tokens = tokenizer.convert_ids_to_tokens(example['input_ids'])

    return {
        "message": "Ejemplo de dato tokenizado",
        "indice": idx,
        "texto_decodificado": text,
        "tokens": tokens[:10],
        "input_ids": example['input_ids'][:5],
        "attention_mask": example['attention_mask'][:5],
        "etiqueta": example['labels']
    }
\end{lstlisting}

% =========================
% Página 71
% =========================

\begin{lstlisting}[style=jsoncode]
{
    "message": "Ejemplo de dato tokenizado",
    "índice": 3701,
    "texto decodificado": "Sinceramente... Después de ver el proceso creativo del logo de totalmente horrible lmao",
    "tokens": [
        "[CLS]",
        "Sinceramente",
        ".",
        ".",
        ".",
        "Después",
        "de",
        "ver",
        "el",
        "proceso"
    ],
    "input_ids": [
        4,
        27235,
        1009,
        1009,
        1009
    ],
    "attention_mask": [
        1,
        1,
        1,
        1,
        1
    ],
    "etiqueta": 0
}
\end{lstlisting}

\section*{A.3. Funciones para el Modelo}

A continuación, se muestran funciones adicionales del proyecto para con el modelo entrenado:

\begin{center}
71
\end{center}

\newpage

% =========================
% Página 72
% =========================

\subsection*{A.3.1. Entrenamiento del Modelo}

\begin{center}
\textbf{Listing A.3. Función para Entrenar el Modelo}
\end{center}

\begin{lstlisting}[style=pythoncode]
def train_model():
    """
    Ejecuta el entrenamiento
    """

    trainer=config_trainer()

    trainer.train()

    best_model_path = Path(os.environ.get("TRANSFORMER_PATH")) / "best_model"
    trainer.save_model(best_model_path)

    return {
        "message": "Modelo entrenado y guardado con éxito.",
        "ruta del modelo": str(best_model_path)
    }
\end{lstlisting}

\subsection*{A.3.2. Métricas del Modelo}

\begin{center}
\textbf{Listing A.4. Función para Obtener Métricas del Modelo}
\end{center}

\begin{lstlisting}[style=pythoncode]
def load_metrics(dataset_type="test"):
    """
    Carga las métricas guardadas desde disco.
    """

    metrics_path = Path(os.environ.get("TRANSFORMER_PATH")) / f"metrics_{dataset_type}.json"

    if not metrics_path.exists():
        raise FileNotFoundError(f"No se encontró el archivo de métricas en: {metrics_path}")

    with open(metrics_path, "r") as f:
        metrics = json.load(f)

    return {
        "message": f"Métricas cargadas desde {dataset_type}.",
        "métricas": metrics
    }
\end{lstlisting}

\begin{center}
72
\end{center}

\newpage

% =========================
% Página 73
% =========================

\subsection*{A.3.3. Pipeline del Modelo}

\begin{center}
\textbf{Listing A.5. Función para Guardar un Pipeline del Modelo}
\end{center}

\begin{lstlisting}[style=pythoncode]
def save_pipeline():
    """
    Guarda el pipeline listo para producción (modelo + tokenizador)
    """

    model_path = Path(os.environ.get("TRANSFORMER_PATH")) / "best_model"
    tokenizer = tok.load_tokenizer()

    # Crear pipeline
    classifier = pipeline(
        "text-classification",
        model=str(model_path),
        tokenizer=tokenizer
    )

    # Guardar pipeline
    pipeline_path = Path(os.environ.get("TRANSFORMER_PATH")) / "pipeline"
    classifier.save_pretrained(pipeline_path)

    return {
        "message": "Pipeline guardado con éxito.",
        "ruta del pipeline": str(pipeline_path)
    }
\end{lstlisting}

\section*{A.4. Ejemplos de textos clasificados}

A continuación, se muestran salidas de textos clasificados por la función \textit{classifier} vista en el capítulo 9.

\begin{lstlisting}[style=jsoncode]
{
    "mensaje": "Clasificación completada.",
    "resultado": {
        "texto": "Impresora negra",
        "etiqueta": "no ofensivo",
        "confianza": 0.9693416357040405
    }
}
\end{lstlisting}

\begin{center}
73
\end{center}

\newpage

% =========================
% Página 74
% =========================

\begin{lstlisting}[style=jsoncode]
{
    "mensaje": "Clasificación completada.",
    "resultado": {
        "texto": "calla negra",
        "etiqueta": "ofensivo",
        "confianza": 0.9984779953956604
    }
}
\end{lstlisting}

\begin{center}
74
\end{center}

\newpage

% =========================
% Página 75
% =========================
\thispagestyle{plain}


